\section*{ÖZET}
\addcontentsline{toc}{section}{Özet}

Bu proje, yaya geçitlerindeki trafik ihlallerini gerçek zamanlı olarak tespit eden, Edge-AI tabanlı akıllı bir sistemin tasarım ve uygulamasını kapsamaktadır. ESP32-S3 mikrodenetleyici üzerinde çalışan kamera modülü, WiFi aracılığıyla MJPEG video akışı sağlamaktadır. Python tabanlı işleme motoru, YOLOv8n nesne tespit modeli kullanarak kişi ve araç tespiti yapmakta; özel geliştirilmiş Canny kenar tespiti ve Hough dönüşümü algoritması ile yaya geçidi bölgesi otomatik olarak belirlenmektedir. İhlal tespitinde anlık fotoğraf kaydı, Telegram bildirimi ve Firebase loglama özellikleri bulunmaktadır. Sistem, düşük maliyetli donanım ile yüksek doğruluklu tespit sağlayan hibrit bir mimari sunmaktadır.

\vspace{1cm}
\noindent\textbf{Anahtar Kelimeler:} Yaya Geçidi Güvenliği, Edge-AI, YOLOv8, ESP32, Nesne Tespiti, Bilgisayarlı Görü

\newpage

\section*{ABSTRACT}
\addcontentsline{toc}{section}{Abstract}

This project covers the design and implementation of an intelligent, Edge-AI-based system that detects traffic violations at pedestrian crossings in real time. The camera module running on an ESP32-S3 microcontroller provides MJPEG video streaming via WiFi. The Python-based processing engine performs person and vehicle detection using the YOLOv8n object detection model, while the custom-developed Canny edge detection and Hough transform algorithm automatically determines the crosswalk region. Violation detection features include instant photo capture, Telegram notification, and Firebase logging. The system presents a hybrid architecture that delivers high-accuracy detection with low-cost hardware.

\vspace{1cm}
\noindent\textbf{Keywords:} Pedestrian Crossing Safety, Edge-AI, YOLOv8, ESP32, Object Detection, Computer Vision
